\section{Stationary-three full C from a complete diagonal entrance}
This theorem treats actual X=1^alpha2^u3^r with
\[
                      4<=u<=32,\quad r>=2u,\quad alpha>=1.
\]
Every current rank and every C shift>=3 is included. The three-count
r and initial padding alpha are UNBOUNDED. The finite binary-count
range is explicit; no theorem for all binary counts is asserted.
The new finite premise is a COMPLETE diagonal entrance, not a random
sample: every u4..32, r0=2u and every needed padding is reconstructed.

\paragraph{A stronger derivative budget, with its hereditary invariant.}
For fixedu>=4, define P(u,r) to mean
\[
 D_n²-Dprev Dnext>=D_n\quad(r+3<=n<ell),                 \tag{1}
\]
where ell is the FIRST actual D mode, and all quantities use the
actual finite arm and its clipped positive derivative. We prove
P(u,r) implies P(u,r+1) for EVERYr>=8 and EVERYpaddingalpha.
No C premise is used in this strong-minor transfer.

Let R_n be the positive derivative of H3=(1-z-z²-z³)^(-1).
For u>=4 the finite binary arm agrees with H2 through degree5,
for every alpha>=1. Consequently H3-F_binary has first THREE
coefficients (1,2,5) at weights3,4,5. Multiplication by S3^r gives
the first three ordinary losses of X relative to H3, at r+3,r+4,r+5:
\[
                       1,\quad r+2,\quad(r²+5r+10)/2.
\]
The first three D losses are therefore
\[
                       1,\quad b=r+1,\quad c=(r²+3r+6)/2. \tag{2}
\]
The finite initial-one correction has NOT been discarded: at u>=4
its first possible relevant ordinary loss is at least r+6, so it
does not change precisely these THREE reads, for ANYalpha.

\paragraph{The two new left entrance centres are analytic.}
The tribonacci adjacent minor M_n=R_n²-Rprev Rnext satisfies the
exterior-square recurrence M_n=-Mprev-Mprevprev+Mprevprevprev.
The three base values at actual centres5,6,7 are3,1,-7.
Absolute-value induction against the positive tribonacci recurrence
gives |M_n|<=Rprev for EVERYn>=5. The full arbitrary-natural
recurrence and domination induction have Lean proofs.

For first current n=r+3, write referenceprevT,currentU,nextV.
Its actual triple is (T,U-1,V-r-1). By U<=2T and Mref>=-T,
\[
 Mactual-(U-1)>= (r-6)T+2>=0\qquad(r>=6).               \tag{3}
\]
For next centre n=r+4 the triple is (U-1,V-b,W-c).
If a clipped successor is zero the centre cannot be strictly before
the FIRSTDmode; in the required band W>=c. Using V<=2U and
Mref>=-V gives
\[
 Mactual-(V-b)>=U(c-4b-4)+b²+b>=0\qquad(r>=7),          \tag{4}
\]
since c-4b-4=(r²-5r-10)/2>=0. Thus every new r has BOTH
left entrance centres paid directly, without an old-tail read there.

\paragraph{Every later pre-mode centre inherits the budget.}
Let the OLD three-count be r and NEW count r+1. The two centres
just paid are r+4,r+5. Consider any later new centre n>=r+6 strictly
before its FIRSTDmode e. Denote the old FIRSTDmode by ell and its
last positive D rank, the FIRST ordinary mode, by R.
The registered signed ordinary splice says new D agrees with
z²[2]+S3Dold through R+1, and is nonincreasing after that rank.
Hence e<=R+1 and EVERYtriple used at n<e agrees with this raw
positive convolution. Its source z²[2] vanishes at these large ranks.
The first window mode is no later than ell+3, so e<=ell+3.

If n<=ell, all THREE sampled old minor centres n-3,n-2,n-1
are at least r+3 and strictly before ell. The old budget (1) applies.
All five old coefficient inputs are positive: their indices are
large and at most the old positive support endpoint. For values
a,b,c,d,e0 the exact convolution minor is
\[
 (b²-ac)+(d²-ce0)+(bc-ad)+(cd-be0)+(bd-ae0).
\]
Old LC makes the last two cross minors nonnegative. Moreover
b(bc-ad-c)=c(b²-ac-b)+a(c²-bd)>=0,
so b>0 gives bc-ad>=c. Thus the convolution minor is at least
b+c+d, the new central value. The POSITIVE-input guard is essential;
a delta convolved with[3] would refute a zero-input version.

If n>=ell+1, then n<=ell+2 because n<e<=ell+3. The exact
second difference of the new raw window is
(Dold_n-Dold_(n-1))-(Dold_(n-3)-Dold_(n-4)).
The entering pair is after the OLD first mode and is nonincreasing.
The outgoing pair is strictly before that mode and its current index
is at least r+3>=11. The actual sole possible rising-side flat is
at finalcapacity+1=4, so the outgoing INTEGER gain is at least one.
The new window therefore has second difference at most -1.
Its own INTEGER gain is at least one too, since n is before the
new first mode and n!=4. Writing A,B,C for its triple,
M=(B-A)²+A(2B-A-C)>=B. This proves (1) at every remaining centre.
Together with (3)--(4), this is the full unbounded r-append proof.

\paragraph{The complete finite diagonal premise and padding reduction.}
For each fixedu>=4 choose r0=2u. The sharpened derivative centre
theorem gives ell<=B=floor((3u+4r0)/2). The first initial-one loss
in the signed derivative is at r0+u+alpha+1. For
\[
                     alpha_* = r0+floor(u/2)
\]
that loss starts at B+1. Thus every larger padding agrees with
alpha_* at ALL derivative triples strictly before its first mode;
no finite-padding extrapolation based only on experiments is used.
It suffices to check exactly1<=alpha<=alpha_*.

For EVERY4<=u<=32 and r0=2u this is EXACTLY1298hosts and93831
strong-minor comparisons. Their34977distinct nonnegative INTEGER
margins M-D are stored in117chunks and checked by ordinary Lean
kernel reduction, with one aggregate report. Independent positive
suffix-box reconstruction must reproduce every host/mode/margin,
the entire parameter grammar and all source hashes. This is the
ONLYnew finite entrance premise. The arbitrary-r tail closure above
then proves (1) for EVERYr>=2u, every alpha, in this binary range.

\paragraph{Strong ordinary minors transfer COMPLETE renewal C.}
At all needed ranks through the actual FIRSTDmode, the proved full
second-source identity gives Dactual<=R, where R is H3's positive
derivative. This is a derivative domination statement, not inferred
just from ordinary coefficientwise domination.
If A is an actual central D value, the relative-LC determinant is
\[
 R_n² Mactual-A² Mref.
\]
Since Mactual>=A, Mref<=R_n and0<=A<=R_n, it is nonnegative.
Before the first ordinary loss at r+3, actual D agrees with R;
the adjacent first-loss boundary only lowers its successor and
improves relative LC. Thus the WHOLE normalized density D/R
is log-concave through the FIRSTDmode, with D2=R2=1 and
D3/R3=1. The existing normalized tilt/twice-cumulative LR theorem
inherits COMPLETE H3 renewal C at ALL shifts>=3. After the first
Dmode every derivative pair is nonincreasing, so C is immediate.
This proves the displayed full-C theorem without any actual forward
CDF LC assertion or substitution for original coefficient UI.

\paragraph{An unrestricted large last jump over the stationary-three block.}
For the same u4..32,r>=2u, every actual additional last capacity L
with L>=3 and
\[
                    L+4>=floor((3u+4r)/2)
\]
has COMPLETE C at all ranks and shifts. Apply the registered
conditional-prefix/old-mode theorem to the newly proved full-C
block. This permits unbounded r,alpha,L and is independent of
initial-padding length. It is a sufficient long-jump condition,
not full C for every last capacity over a three-block prefix.

Ordinary-product coefficient-tilt UI against ANYopposite and the
registered BOTH-Frising every-depth original C-to-S resource follow,
with old head/physical/current hypotheses retained. General C
outside these families, full weighted A2 and the ORIGINAL general
unamplified one-descent coefficient UI remain OPEN.
